% =====================================================================
% APPENDIX B — DATABASE STRUCTURE AND CODE REPOSITORY
% =====================================================================

\chapter{Database Structure and Code Repository}
\label{app:data}

\section*{B.1. SQLite Database Structure}

The database \verb|kz_energy_balance.db| (size 2.7 MB) contains 18 tables
with a total of 34{,}362 rows.

\begin{table}[h!]
\centering
\caption{List of tables in the consolidated database}
\label{tab:db_tables}
\small
\begin{tabular}{@{}llrr@{}}
\toprule
\# & Table & Rows & Columns \\
\midrule
1 & balance\_long & 8{,}277 & 7 \\
2 & balance\_wide\_by\_fuel & 332 & 134 \\
3 & by\_fuel\_group\_TJ & 132 & 4 \\
4 & electricity\_balance & 8 & 92 \\
5 & kegoc\_balance & 18 & 7 \\
6 & kegoc\_balance\_wide & 9 & 3 \\
7 & korem\_hourly & 21{,}792 & 10 \\
8 & korem\_monthly & 30 & 8 \\
9 & meta & 5 & 2 \\
10 & mirror\_test\_dm & 18 & 6 \\
11 & mirror\_test\_metrics & 12 & 8 \\
12 & ml\_diebold\_mariano & 4 & 5 \\
13 & ml\_dm\_renewable\_cf & 8 & 6 \\
14 & ml\_metrics & 5 & 6 \\
15 & ml\_metrics\_all\_targets & 17 & 7 \\
16 & ml\_metrics\_renewable\_cf & 10 & 7 \\
17 & ml\_predictions & 45 & 7 \\
18 & ml\_predictions\_renewable\_cf & 3{,}640 & 5 \\
\midrule
& \textbf{Total} & \textbf{34{,}362} & --- \\
\bottomrule
\end{tabular}
\end{table}

\section*{B.2. Repository Structure}

The full project repository (relative to its root) is organized as follows:

{\footnotesize
\begin{verbatim}
Dissertation/
|-- Dissertation_Document/        # This thesis (LaTeX)
|   |-- main.tex                  # Entry point
|   |-- refs.bib                  # Bibliography
|   |-- chapters/                 # All chapters
|   |-- figures/                  # Figures
|   `-- compile.bat               # Compile script
|
|-- 01_AI_DataCenter_Kazakhstan/  # Article 1: AI Data Centers
|
|-- 02_Energy_Balance_and_        # Article 2: Energy Balance
|       Demand_Forecasting_KZ/
|   `-- Article/
|       |-- Final/
|       |   |-- main.tex          # Article 2 manuscript
|       |   |-- Energy.pbix       # Power BI dashboard
|       |   `-- figures/          # Figures
|       |-- data/clean/
|       |   `-- kz_energy_balance.db
|       `-- scripts/              # ETL + forecasting (Python)
|           |-- forecast_renewable_cf.py
|           `-- transfer_forecasting/
|
`-- 03_ML_Renewable_Forecasting_  # Exploratory CNN-LSTM draft
        CNN_LSTM/                 # (not used in this thesis)
\end{verbatim}
}

\section*{B.3. Reproducibility}

For full reproducibility of the analysis, execute (paths relative to project root):

{\footnotesize
\begin{verbatim}
# 1. Install dependencies
cd 02_Energy_Balance_and_Demand_Forecasting_KZ
pip install -r requirements.txt

# 2. Run the ETL pipeline (data parser)
cd Article
python scripts/data_parser.py

# 3. Wind capacity-factor forecasting pipeline
#    (6-model comparison + SARIMA-LSTM hybrid)
python scripts/forecast_renewable_cf.py

# 4. Inspect the database
python ../db_info.py

# 5. Open the Power BI dashboard
start Final/Energy.pbix      # Windows
open  Final/Energy.pbix      # macOS

# 6. Compile the manuscript (pdflatex x3 + bibtex)
cd Final
pdflatex main.tex
bibtex main
pdflatex main.tex
pdflatex main.tex

# 7. Compile this thesis document
cd ../../../Dissertation_Document
compile.bat                  # Windows
\end{verbatim}
}

\section*{B.4. Software and Library Versions}

\begin{itemize}
    \item Python 3.11+
    \item pandas 2.x, numpy 1.26+
    \item statsmodels 0.14+
    \item tensorflow 2.15+, keras 2.15+
    \item scikit-learn 1.3+
    \item matplotlib 3.8+, seaborn 0.13+
    \item SQLite 3.40+
    \item Power BI Desktop 2.153+
    \item MiKTeX 25.x or TeX Live 2024+
\end{itemize}
